Write the treatment rows of \(A_{\mathrm b}^{\dagger}\), in the column order \((a_{\mathrm b},z_{\mathrm b})\), as
\[
w_0=\tfrac12e_0\otimes\psi^+,
\qquad
w_1=\tfrac12e_1\otimes\psi^+,
\qquad
w_2=\tfrac12\mathbf 1\otimes\psi^-.
\]
These are exactly the rows prescribed by \(\texttt{def:binary-bow-menu}\) and \(\texttt{def:bow-probe-matrix}\).  The treatment effects obey
\[
\tfrac12e_0+\tfrac12e_1+\tfrac12\mathbf 1=\mathbf 1,
\]
and the terminal effects obey \(e_0+e_1=\mathbf 1\), proving \(\texttt{ass:instrument-normalization}\) at both vertices.  The effects \(\tfrac12e_0,\tfrac12e_1\) span \(\mathbb R^2\).  Moreover,
\[
\det\begin{pmatrix}3/4&1/4\\1/4&3/4\end{pmatrix}=\frac12\neq0,
\]
so \(\psi^+,\psi^-\) span \(\mathbb R^2\).  At the terminal vertex, both the effects and states are the coordinate basis \(e_0,e_1\), and for \(i\in\{0,1\}\),
\[
(e_i\otimes e_i)(a_{\mathrm t},z_{\mathrm t})
=\mathbf 1\{a_{\mathrm t}=i,z_{\mathrm t}=i\}
=d_i(a_{\mathrm t},z_{\mathrm t}).
\]

If \(\alpha w_0+\beta w_1+\gamma w_2=0\), restriction to the coordinates with \(a_{\mathrm b}=0\) gives \(\alpha\psi^++\gamma\psi^-=0\).  The nonzero determinant gives \(\alpha=\gamma=0\).  Restriction to \(a_{\mathrm b}=1\) then gives \(\beta\psi^+=0\), hence \(\beta=0\).  Therefore
\[
\operatorname{rank}(A_{\mathrm b}^{\dagger})=3<4.
\]
Direct coordinate calculation gives
\[
e_0=\tfrac32\psi^+-\tfrac12\psi^-,
\qquad
e_1=-\tfrac12\psi^++\tfrac32\psi^-.
\]
Since \(r_i=\mathbf 1\otimes e_i\),
\[
r_0=3w_0+3w_1-w_2,
\qquad
r_1=-w_0-w_1+3w_2.
\]
Thus both target rows lie in \(\operatorname{row}(A_{\mathrm b}^{\dagger})\), and every \(e_y\), \(y\in\{0,1\}\), lies in the terminal effect span.

Fix any \(\mathfrak C\) between the primitive-positive and full finite edge-source classes.  Applying the success direction of \(\texttt{thm:bow-rowspan-iff}\) to each \((z^{\circ},y)\in\{0,1\}^2\) proves \(\operatorname{GID}_{\mathcal I_{\mathrm b}^{\dagger}}^{\mathfrak C}(Q_{z^{\circ},y})\) for all four atoms.  In the identifying functional one may take treatment coefficient vectors \((3,3,-1)\) for \(z^{\circ}=0\) and \((-1,-1,3)\) for \(z^{\circ}=1\).  Taking \(\mathfrak C=\mathfrak M_{\mathrm{fin}}^{\mathrm{edge}}(G_{\mathrm b},\mathcal A)\) proves the full-class claim.  Since the treatment table has four coordinates but the treatment probe matrix has rank three, the probe map is not injective on that full table.  Identification is therefore query specific and non-tomographic.